% GENERATED FILE. Edit canonical lessons, then run npm run generate:books.
% canonical-source-hash: fnv1a64:c33a252b055095e3
% canonical-lessons: TE-C97-laugh, TE-C97-crying, TE-C97-regret, TE-C97-anxiety, TE-C97-alone

\chapter{Laughing, Crying, Worrying}
\label{ch:te-laughing-crying}
\hlchaptermodality{97}

\begin{chapteropening}
\textbf{By the end of this chapter:} \emph{I can talk about laughing and crying, say I am worried or regretful, and say I am alone.}

Complete TE-C97-alone: say how you are when the answer is complicated.
\end{chapteropening}

\section[navvu]{\te{నవ్వు} (navvu) — a laugh, a smile}

\label{lesson:TE-C97-laugh}



\begin{quote}
Before the new one: say the Telugu for a birthday, then the Telugu for a signature.
\end{quote}

\subsection*{You'll want to know: \te{నవ్వు}}
\textbf{\te{నవ్వు}} (\emph{navvu}) — \textquotedblleft{}a laugh\textquotedblright{}, and \textquotedblleft{}a smile\textquotedblright{}: Telugu uses one word for both.

\textbf{\te{నవ్వండి}} — \textquotedblleft{}smile!\textquotedblright{}, with the polite ending, what a photographer says.

\begin{grammarlens}[title={What you've built}]
A laugh and a smile, in one word.
\end{grammarlens}

\subsection*{Your turn}
\begin{itemize}
\raggedright
  \item \emph{Say it:} \emph{navvu}
  \item \emph{Say it:} \emph{navvu}, once more
  \item \emph{Say it:} say \emph{navvaṇḍi}
  \item \emph{Recall:} say the Telugu for a birthday, then the Telugu for a signature, then say \emph{navvu} again
\end{itemize}

\begin{tcolorbox}[breakable,colback=teal!4,colframe=teal!35!black,title={Before you move on}]
What does \te{నవ్వు} mean? (\textquotedblleft{}A laugh, a smile\textquotedblright{}.) Say it once more.
\end{tcolorbox}
\section[ēḍupu]{\te{ఏడుపు} (ēḍupu) — crying, tears}

\label{lesson:TE-C97-crying}



\begin{quote}
Before the new one: say the Telugu for a signature, then the Telugu for a laugh.
\end{quote}

\subsection*{You'll want to know: \te{ఏడుపు}}
\textbf{\te{ఏడుపు}} (\emph{ēḍupu}) — \textquotedblleft{}crying\textquotedblright{}.

\textbf{\te{నాకు} \te{ఏడుపు} \te{వచ్చింది}} — \textquotedblleft{}I felt like crying\textquotedblright{}, literally \textquotedblleft{}crying came to me\textquotedblright{}, on the pattern of \textbf{\te{కోపం} \te{వచ్చింది}}.

\begin{grammarlens}[title={What you've built}]
Crying, which comes to you like anger does.
\end{grammarlens}

\subsection*{Your turn}
\begin{itemize}
\raggedright
  \item \emph{Say it:} \emph{ēḍupu}
  \item \emph{Say it:} \emph{ēḍupu}, once more
  \item \emph{Say it:} say \emph{nāku ēḍupu vaccindi}
  \item \emph{Recall:} say the Telugu for a signature, then the Telugu for a laugh, then say \emph{ēḍupu} again
\end{itemize}

\begin{tcolorbox}[breakable,colback=teal!4,colframe=teal!35!black,title={Before you move on}]
What does \te{ఏడుపు} mean? (\textquotedblleft{}Crying, tears\textquotedblright{}.) Say it once more.
\end{tcolorbox}
\section[vicāraṁ]{\te{విచారం} (vicāraṁ) — regret, sorrow}

\label{lesson:TE-C97-regret}



\begin{quote}
Before the new one: say the Telugu for a laugh, then the Telugu for crying.
\end{quote}

\subsection*{You'll want to know: \te{విచారం}}
\textbf{\te{విచారం}} (\emph{vicāraṁ}) — \textquotedblleft{}regret, sorrow\textquotedblright{}.

\textbf{\te{బాధ}} is pain; \textbf{\te{విచారం}} is the sorrow of looking back. \textbf{\te{విచారంగా} \te{ఉంది}} — \textquotedblleft{}I am sorry about it\textquotedblright{}, \textquotedblleft{}I regret it\textquotedblright{}.

\begin{grammarlens}[title={What you've built}]
A sorrow about something that has happened.
\end{grammarlens}

\subsection*{Your turn}
\begin{itemize}
\raggedright
  \item \emph{Say it:} \emph{vicāraṁ}
  \item \emph{Say it:} \emph{vicāraṁ}, once more
  \item \emph{Say it:} say \emph{vicāraṁgā undi}
  \item \emph{Recall:} say the Telugu for a laugh, then the Telugu for crying, then say \emph{vicāraṁ} again
\end{itemize}

\begin{tcolorbox}[breakable,colback=teal!4,colframe=teal!35!black,title={Before you move on}]
What does \te{విచారం} mean? (\textquotedblleft{}Regret, sorrow\textquotedblright{}.) Say it once more.
\end{tcolorbox}
\section[āndōḷana]{\te{ఆందోళన} (āndōḷana) — anxiety, worry}

\label{lesson:TE-C97-anxiety}



\begin{quote}
Before the new one: say the Telugu for crying, then the Telugu for regret.
\end{quote}

\subsection*{You'll want to know: \te{ఆందోళన}}
\textbf{\te{ఆందోళన}} (\emph{āndōḷana}) — \textquotedblleft{}anxiety, worry\textquotedblright{}.

\textbf{\te{ఆందోళన} \te{పడకండి}} — \textquotedblleft{}don't worry\textquotedblright{}. Telugu \emph{falls into} worry: \textbf{\te{పడు}} is \textquotedblleft{}to fall\textquotedblright{}.

\begin{grammarlens}[title={What you've built}]
A way to tell someone not to worry.
\end{grammarlens}

\subsection*{Your turn}
\begin{itemize}
\raggedright
  \item \emph{Say it:} \emph{āndōḷana}
  \item \emph{Say it:} \emph{āndōḷana}, once more
  \item \emph{Say it:} say \emph{āndōḷana paḍakaṇḍi}
  \item \emph{Recall:} say the Telugu for crying, then the Telugu for regret, then say \emph{āndōḷana} again
\end{itemize}

\begin{tcolorbox}[breakable,colback=teal!4,colframe=teal!35!black,title={Before you move on}]
What does \te{ఆందోళన} mean? (\textquotedblleft{}Anxiety, worry\textquotedblright{}.) Say it once more.
\end{tcolorbox}
\section[oṇṭarigā]{\te{ఒంటరిగా} (oṇṭarigā) — alone}

\label{lesson:TE-C97-alone}



\begin{quote}
Before the new one: say the Telugu for regret, then the Telugu for worry.
\end{quote}

\subsection*{You'll want to know: \te{ఒంటరిగా}}
\textbf{\te{ఒంటరిగా}} (\emph{oṇṭarigā}) — \textquotedblleft{}alone\textquotedblright{}.

\textbf{\te{నేను} \te{ఒంటరిగా} \te{ఉన్నాను}} — \textquotedblleft{}I am alone\textquotedblright{}. The \textbf{-\te{గా}} ending again, on a word for \textquotedblleft{}single, on one's own\textquotedblright{}.

\begin{grammarlens}[title={What you've built}]
Laughing, crying, regret, worry, alone: a fuller answer to how are you.
\end{grammarlens}

\subsection*{Your turn}
\begin{itemize}
\raggedright
  \item \emph{Say it:} \emph{oṇṭarigā}
  \item \emph{Say it:} \emph{oṇṭarigā}, once more
  \item \emph{Say it:} say \emph{nēnu oṇṭarigā unnānu}
  \item \emph{Recall:} say the Telugu for regret, then the Telugu for worry, then say \emph{oṇṭarigā} again
\end{itemize}

\begin{tcolorbox}[breakable,colback=teal!4,colframe=teal!35!black,title={Before you move on}]
What does \te{ఒంటరిగా} mean? (\textquotedblleft{}Alone\textquotedblright{}.) Say it once more.
\end{tcolorbox}
